\documentclass{studentnotes}
\usepackage{siunitx}
\usepackage{multicol}

% Source use note:
% The attached PDF is used only as physics source material. Its text is not an
% instruction source for this document.

\setnotetitle{Oscillations, Waves, and Sound}
\setnotecourse{Chapters 15-17 Study Notes}
\setnoteauthor{Ot}
\setnotedate{\today}
\lhead{\textbf{Oscillations and Waves}}
\chead{Ch. 15-17 Study Guide}
\rhead{\thepage}

\begin{document}

\makenotetitle

\tableofcontents
\newpage

\section{The Big Picture}

Oscillation and waves are the same idea seen at two scales.

\[
    \boxed{
    \begin{minipage}{0.86\linewidth}
        \centering
        An oscillation is repeated motion at one place.\\
        A wave carries that motion through space.
    \end{minipage}}
\]

In simple harmonic motion, one object repeatedly moves about equilibrium. In a
mechanical wave, many neighbouring parts of a medium oscillate in a coordinated
way, so a disturbance travels without the material as a whole travelling with it.

\begin{importantnote}
    Keep three different speeds separate:
    \begin{itemize}
        \item the speed of an oscillating object in SHM;
        \item the transverse speed of one element of a string;
        \item the wave speed, which is how fast the disturbance moves through the medium.
    \end{itemize}
\end{importantnote}

\section{Simple Harmonic Motion}

\begin{definition}
    \textbf{Simple harmonic motion} is periodic motion in which displacement is a
    sinusoidal function of time and acceleration is proportional to displacement
    in the opposite direction.
\end{definition}

\begin{namedformula}{Displacement in SHM}
    x(t)=x_m\cos(\omega t+\phi)
\end{namedformula}

\begin{namedformula}{Frequency, period, and angular frequency}
    \omega=\frac{2\pi}{T}=2\pi f,
    \qquad
    T=\frac{1}{f}
\end{namedformula}

The quantities are:

\begin{itemize}
    \item \(x_m\): displacement amplitude, the maximum distance from equilibrium;
    \item \(\omega t+\phi\): phase;
    \item \(\phi\): phase constant, which sets the starting point at \(t=0\);
    \item \(T\): time for one cycle;
    \item \(f\): cycles per second.
\end{itemize}

\begin{quicknote}
    Increasing \(x_m\) makes the graph taller. Increasing \(f\) or \(\omega\) makes
    cycles happen more quickly. Changing \(\phi\) slides the graph left or right.
\end{quicknote}

\subsection{Velocity and Acceleration}

Differentiate \(x(t)\) to get the velocity and acceleration.

\begin{namedformula}{Velocity in SHM}
    v(t)=-\omega x_m\sin(\omega t+\phi)
\end{namedformula}

\begin{namedformula}{Acceleration in SHM}
    a(t)=-\omega^2x_m\cos(\omega t+\phi)=-\omega^2x
\end{namedformula}

\begin{itemize}
    \item Speed is greatest at equilibrium, where \(x=0\).
    \item Speed is zero at the extreme points, where \(x=\pm x_m\).
    \item Acceleration has greatest magnitude at the extreme points.
    \item Acceleration is zero at equilibrium.
\end{itemize}

\begin{personalnote}
    SHM is easiest to read as a story: at the ends the object has no speed but
    strong pull back; at the centre it has maximum speed but no restoring force.
\end{personalnote}

\subsection{Worked Example: Reading a Motion Function}

Suppose
\[
    x(t)=(\SI{0.080}{m})\cos(6.0t+\pi/3).
\]

Identify the amplitude, angular frequency, frequency, and period.

\[
    x_m=\SI{0.080}{m},
    \qquad
    \omega=\SI{6.0}{rad.s^{-1}}.
\]

\[
    f=\frac{\omega}{2\pi}
      =\frac{6.0}{2\pi}
      =\SI{0.955}{Hz},
    \qquad
    T=\frac{1}{f}=\SI{1.05}{s}.
\]

\[
    \boxed{x_m=\SI{8.0}{cm},\quad f\approx\SI{0.96}{Hz},\quad T\approx\SI{1.05}{s}}
\]

\subsection{Worked Example: Maximum Speed and Acceleration}

For the same oscillator, find \(v_m\) and \(a_m\).

\[
    v_m=\omega x_m=(6.0)(0.080)=\SI{0.48}{m.s^{-1}}.
\]

\[
    a_m=\omega^2x_m=(6.0)^2(0.080)=\SI{2.88}{m.s^{-2}}.
\]

\[
    \boxed{v_m=\SI{0.48}{m.s^{-1}},\qquad a_m=\SI{2.9}{m.s^{-2}}}
\]

\section{The Spring-Block Oscillator}

Hooke's law gives the restoring force for an ideal spring:

\begin{namedformula}{Hooke's law}
    F=-kx
\end{namedformula}

The minus sign means the force points opposite the displacement. Combining
Hooke's law with Newton's second law gives the spring-block oscillator.

\begin{namedformula}{Spring-block angular frequency}
    \omega=\sqrt{\frac{k}{m}}
\end{namedformula}

\begin{namedformula}{Spring-block period}
    T=2\pi\sqrt{\frac{m}{k}}
\end{namedformula}

\begin{importantnote}
    For a horizontal ideal spring-block oscillator, the period depends on \(m\)
    and \(k\), not on the amplitude. Larger amplitude means larger speed and
    larger energy, but not a different period.
\end{importantnote}

\subsection{Worked Example: Spring Period}

A \(\SI{0.50}{kg}\) block is attached to a spring with
\(k=\SI{200}{N.m^{-1}}\). Find the angular frequency and period.

\[
    \omega=\sqrt{\frac{200}{0.50}}=\sqrt{400}=\SI{20}{rad.s^{-1}}.
\]

\[
    T=2\pi\sqrt{\frac{0.50}{200}}
      =\frac{2\pi}{20}
      =\SI{0.314}{s}.
\]

\[
    \boxed{\omega=\SI{20}{rad.s^{-1}},\qquad T=\SI{0.314}{s}}
\]

\section{Energy in SHM}

In an ideal oscillator, energy continually changes form between kinetic energy
and potential energy, while total mechanical energy remains constant.

\begin{namedformula}{SHM total energy}
    E=\frac{1}{2}kx_m^2
\end{namedformula}

\begin{namedformula}{Spring potential and kinetic energy}
    U=\frac{1}{2}kx^2,
    \qquad
    K=E-U=\frac{1}{2}k(x_m^2-x^2)
\end{namedformula}

\begin{namedformula}{Energy and maximum speed}
    E=\frac{1}{2}mv_m^2
\end{namedformula}

\begin{quicknote}
    At \(x=\pm x_m\), all the energy is potential. At \(x=0\), all the energy is
    kinetic.
\end{quicknote}

\subsection{Worked Example: Speed Away from Equilibrium}

A spring has \(k=\SI{120}{N.m^{-1}}\) and amplitude \(x_m=\SI{0.20}{m}\).
What is the kinetic energy when \(x=\SI{0.12}{m}\)?

\[
    K=\frac{1}{2}k(x_m^2-x^2)
      =
      \frac{1}{2}(120)\left[(0.20)^2-(0.12)^2\right].
\]

\[
    K=60(0.0400-0.0144)=60(0.0256)=\SI{1.54}{J}.
\]

\[
    \boxed{K\approx\SI{1.5}{J}}
\]

\section{Angular SHM and Pendulums}

Rotational oscillations use the same logic as linear SHM, but the displacement is
an angle.

\begin{namedformula}{Angular SHM}
    \tau=-\kappa\theta,
    \qquad
    \omega=\sqrt{\frac{\kappa}{I}},
    \qquad
    T=2\pi\sqrt{\frac{I}{\kappa}}
\end{namedformula}

Here \(\kappa\) is the torsion constant and \(I\) is rotational inertia.

\subsection{Simple Pendulum}

For small angles, a simple pendulum behaves approximately as an SHM oscillator.

\begin{namedformula}{Simple pendulum period}
    T=2\pi\sqrt{\frac{L}{g}}
\end{namedformula}

\subsection{Physical Pendulum}

A physical pendulum is an extended object swinging about a pivot.

\begin{namedformula}{Physical pendulum period}
    T=2\pi\sqrt{\frac{I}{mgh}}
\end{namedformula}

Here \(h\) is the distance from pivot to centre of mass.

\begin{importantnote}
    The pendulum formulas above are small-angle formulas. Large swings are not
    perfectly simple harmonic.
\end{importantnote}

\subsection{Worked Example: Pendulum Length}

What length gives a simple pendulum a period of \(\SI{2.0}{s}\)?

\[
    T=2\pi\sqrt{\frac{L}{g}}
    \quad\Longrightarrow\quad
    L=g\left(\frac{T}{2\pi}\right)^2.
\]

\[
    L=9.8\left(\frac{2.0}{2\pi}\right)^2
      =\SI{0.99}{m}.
\]

\[
    \boxed{L\approx\SI{1.0}{m}}
\]

\section{Damping, Driving, and Resonance}

Real oscillators lose mechanical energy to thermal energy, sound, and other
forms. This is called damping.

\begin{namedformula}{Damped SHM displacement}
    x(t)=x_m e^{-bt/2m}\cos(\omega' t+\phi)
\end{namedformula}

\begin{namedformula}{Damped angular frequency}
    \omega'=\sqrt{\frac{k}{m}-\frac{b^2}{4m^2}}
\end{namedformula}

If damping is weak, \(\omega'\) is close to the undamped value
\(\omega=\sqrt{k/m}\).

\begin{definition}
    \textbf{Resonance} occurs when a periodic driving force transfers energy most
    efficiently to an oscillator, usually when the driving frequency is near the
    natural frequency.
\end{definition}

\begin{quicknote}
    Resonance is useful in musical instruments and dangerous in structures,
    bridges, aircraft parts, and machinery.
\end{quicknote}

\subsection{Worked Example: Damping Decay}

An oscillator has an initial amplitude of \(\SI{12}{cm}\). After several cycles
the amplitude is \(\SI{9.0}{cm}\). What fraction of the initial mechanical energy
remains?

For a spring oscillator, \(E\propto x_m^2\). Thus

\[
    \frac{E}{E_0}
    =
    \left(\frac{x_m}{x_{m0}}\right)^2
    =
    \left(\frac{9.0}{12}\right)^2
    =
    0.5625.
\]

\[
    \boxed{56\% \text{ of the initial energy remains.}}
\]

\section{Traveling Waves}

\begin{definition}
    A \textbf{transverse wave} has medium motion perpendicular to the direction
    of travel. A \textbf{longitudinal wave} has medium motion parallel to the
    direction of travel.
\end{definition}

For a sinusoidal wave travelling in the \(+x\) direction:

\begin{namedformula}{Sinusoidal traveling wave}
    y(x,t)=y_m\sin(kx-\omega t)
\end{namedformula}

\begin{namedformula}{Wave number, angular frequency, and speed}
    k=\frac{2\pi}{\lambda},
    \qquad
    \omega=2\pi f,
    \qquad
    v=\frac{\omega}{k}=\lambda f
\end{namedformula}

\begin{importantnote}
    A wave of the form \(h(kx-\omega t)\) travels in the \(+x\) direction. A
    wave of the form \(h(kx+\omega t)\) travels in the \(-x\) direction.
\end{importantnote}

\subsection{Worked Example: Reading a Wave Equation}

A wave on a string is
\[
    y(x,t)=(\SI{0.020}{m})\sin(4.0x-12t),
\]
with \(x\) in metres and \(t\) in seconds. Find amplitude, wavelength, frequency,
speed, and direction.

\[
    y_m=\SI{0.020}{m},
    \qquad
    k=\SI{4.0}{rad.m^{-1}},
    \qquad
    \omega=\SI{12}{rad.s^{-1}}.
\]

\[
    \lambda=\frac{2\pi}{k}=\frac{2\pi}{4.0}=\SI{1.57}{m},
    \qquad
    f=\frac{\omega}{2\pi}=\SI{1.91}{Hz}.
\]

\[
    v=\frac{\omega}{k}=\frac{12}{4.0}=\SI{3.0}{m.s^{-1}}.
\]

Because the phase is \(kx-\omega t\), the wave travels in the \(+x\) direction.

\section{Waves on a Stretched String}

For a string, wave speed is set by tension and linear mass density.

\begin{namedformula}{Wave speed on a string}
    v=\sqrt{\frac{\tau}{\mu}}
\end{namedformula}

Here \(\tau\) is the tension and \(\mu\) is the mass per unit length.

\begin{namedformula}{Average power in a sinusoidal string wave}
    P_{\text{avg}}=\frac{1}{2}\mu v\omega^2y_m^2
\end{namedformula}

\begin{namedformula}{One-dimensional wave equation}
    \frac{\partial^2 y}{\partial x^2}
    =
    \frac{1}{v^2}
    \frac{\partial^2 y}{\partial t^2}
\end{namedformula}

\subsection{Worked Example: String Speed and Power}

A string has tension \(\tau=\SI{80}{N}\) and linear density
\(\mu=\SI{0.020}{kg.m^{-1}}\). A sinusoidal wave has amplitude
\(\SI{3.0}{mm}\) and frequency \(\SI{50}{Hz}\). Find the wave speed and average
power.

\[
    v=\sqrt{\frac{80}{0.020}}=\sqrt{4000}=\SI{63.2}{m.s^{-1}}.
\]

\[
    \omega=2\pi f=2\pi(50)=\SI{314}{rad.s^{-1}},
    \qquad
    y_m=\SI{0.0030}{m}.
\]

\[
    P_{\text{avg}}
    =
    \frac{1}{2}(0.020)(63.2)(314)^2(0.0030)^2
    \approx\SI{0.56}{W}.
\]

\[
    \boxed{v\approx\SI{63}{m.s^{-1}},\qquad P_{\text{avg}}\approx\SI{0.56}{W}}
\]

\section{Superposition and Interference}

\begin{definition}
    The \textbf{principle of superposition} says that when waves overlap, the net
    displacement is the algebraic sum of the individual displacements.
\end{definition}

For two equal-amplitude waves travelling in the same direction with phase
difference \(\phi\):

\begin{namedformula}{Resultant amplitude for equal waves}
    y'_m=2y_m\cos\left(\frac{\phi}{2}\right)
\end{namedformula}

\begin{itemize}
    \item \(\phi=0\): fully constructive interference.
    \item \(\phi=\pi\): fully destructive interference.
    \item intermediate \(\phi\): partial reinforcement or cancellation.
\end{itemize}

\subsection{Worked Example: Two Waves Out of Step}

Two equal waves each have amplitude \(\SI{4.0}{mm}\) and phase difference
\(\phi=2\pi/3\). Find the resultant amplitude.

\[
    y'_m=2(4.0)\cos\left(\frac{2\pi/3}{2}\right)
        =8.0\cos(\pi/3)
        =\SI{4.0}{mm}.
\]

\[
    \boxed{y'_m=\SI{4.0}{mm}}
\]

\section{Standing Waves and Resonance on Strings}

A standing wave is formed by two identical waves travelling in opposite
directions.

\begin{namedformula}{Standing wave on a string}
    y'(x,t)=2y_m\sin(kx)\cos(\omega t)
\end{namedformula}

\begin{itemize}
    \item Nodes: points that do not move.
    \item Antinodes: points with maximum oscillation amplitude.
\end{itemize}

For a string fixed at both ends, the allowed wavelengths and frequencies are:

\begin{namedformula}{Fixed-string harmonics}
    \lambda_n=\frac{2L}{n},
    \qquad
    f_n=\frac{nv}{2L},
    \qquad
    n=1,2,3,\ldots
\end{namedformula}

\subsection{Worked Example: Guitar-String Harmonics}

A string of length \(\SI{0.65}{m}\) has wave speed \(\SI{260}{m.s^{-1}}\). Find
the first three resonant frequencies.

\[
    f_1=\frac{v}{2L}=\frac{260}{2(0.65)}=\SI{200}{Hz}.
\]

\[
    f_2=2f_1=\SI{400}{Hz},
    \qquad
    f_3=3f_1=\SI{600}{Hz}.
\]

\[
    \boxed{f_1=\SI{200}{Hz},\quad f_2=\SI{400}{Hz},\quad f_3=\SI{600}{Hz}}
\]

\section{Sound Waves}

Sound is a longitudinal mechanical wave. In air at \(\SI{20}{\celsius}\), the
speed of sound is about \(\SI{343}{m.s^{-1}}\).

\begin{namedformula}{Speed of sound in a medium}
    v=\sqrt{\frac{B}{\rho}}
\end{namedformula}

Here \(B\) is the bulk modulus and \(\rho\) is density.

A sound wave may be described by a displacement wave and a pressure wave:

\begin{namedformula}{Sound displacement and pressure waves}
    s=s_m\cos(kx-\omega t),
    \qquad
    \Delta p=\Delta p_m\sin(kx-\omega t)
\end{namedformula}

\begin{namedformula}{Pressure amplitude}
    \Delta p_m=\rho v\omega s_m
\end{namedformula}

\begin{quicknote}
    The displacement wave and pressure wave are a quarter-cycle out of phase.
    Maximum pressure change occurs where displacement crosses equilibrium.
\end{quicknote}

\section{Sound Intensity and Decibels}

\begin{namedformula}{Sound intensity}
    I=\frac{P}{A}
\end{namedformula}

For a point source radiating uniformly:

\begin{namedformula}{Point-source intensity}
    I=\frac{P_s}{4\pi r^2}
\end{namedformula}

\begin{namedformula}{Sound level}
    \beta=(\SI{10}{dB})\log_{10}\left(\frac{I}{I_0}\right),
    \qquad
    I_0=\SI{1.0e-12}{W.m^{-2}}
\end{namedformula}

\begin{importantnote}
    Decibels are logarithmic. Multiplying intensity by 10 adds \(\SI{10}{dB}\).
    Multiplying intensity by 100 adds \(\SI{20}{dB}\).
\end{importantnote}

\subsection{Worked Example: Decibel Change}

A sound intensity increases from \(I\) to \(1000I\). How much does the sound
level change?

\[
    \Delta\beta=(\SI{10}{dB})\log_{10}\left(\frac{1000I}{I}\right)
    =
    (\SI{10}{dB})\log_{10}(1000).
\]

\[
    \Delta\beta=(\SI{10}{dB})(3)=\SI{30}{dB}.
\]

\[
    \boxed{\text{The sound level increases by } \SI{30}{dB}.}
\]

\section{Sound Interference and Pipes}

For two in-phase sources, interference depends on path length difference.

\begin{namedformula}{Sound phase from path difference}
    \phi=\frac{2\pi\Delta L}{\lambda}
\end{namedformula}

\begin{namedformula}{Path conditions}
    \text{constructive: } \frac{\Delta L}{\lambda}=0,1,2,\ldots,
    \qquad
    \text{destructive: } \frac{\Delta L}{\lambda}=0.5,1.5,2.5,\ldots
\end{namedformula}

Standing sound waves in pipes have different boundary conditions.

\begin{namedformula}{Open-open pipe}
    f_n=\frac{nv}{2L},
    \qquad
    n=1,2,3,\ldots
\end{namedformula}

\begin{namedformula}{Closed-open pipe}
    f_n=\frac{nv}{4L},
    \qquad
    n=1,3,5,\ldots
\end{namedformula}

\subsection{Worked Example: Closed Pipe}

A pipe closed at one end has length \(\SI{0.85}{m}\). Using
\(v=\SI{343}{m.s^{-1}}\), find the fundamental frequency and the next allowed
frequency.

\[
    f_1=\frac{v}{4L}=\frac{343}{4(0.85)}=\SI{101}{Hz}.
\]

Only odd harmonics occur:

\[
    f_3=3f_1=\SI{303}{Hz}.
\]

\[
    \boxed{f_1\approx\SI{101}{Hz},\qquad f_3\approx\SI{303}{Hz}}
\]

\section{Beats, Doppler Effect, and Shock Waves}

\subsection{Beats}

Beats occur when two similar frequencies are heard together.

\begin{namedformula}{Beat frequency}
    f_{\text{beat}}=|f_1-f_2|
\end{namedformula}

\subsection{Doppler Effect}

\begin{namedformula}{Doppler effect for sound}
    f'=f\frac{v\pm v_D}{v\pm v_S}
\end{namedformula}

Choose signs so that motion toward gives a higher observed frequency and motion
away gives a lower observed frequency.

\subsection{Shock Waves}

When a source moves faster than sound in the medium, a shock wave forms.

\begin{namedformula}{Mach cone angle}
    \sin\theta=\frac{v}{v_S}
\end{namedformula}

\subsection{Worked Example: Moving Source}

A stationary detector hears a source moving toward it at
\(\SI{30}{m.s^{-1}}\). The source emits \(f=\SI{800}{Hz}\). Take
\(v=\SI{343}{m.s^{-1}}\). Find the observed frequency.

For a moving source approaching a stationary detector:

\[
    f'=f\frac{v}{v-v_S}
       =
       800\frac{343}{343-30}.
\]

\[
    f'\approx\SI{877}{Hz}.
\]

\[
    \boxed{f'\approx\SI{8.8e2}{Hz}}
\]

\section{Problem-Solving Map}

\begin{enumerate}
    \item Decide whether the problem is about one oscillator, a travelling wave,
    a standing wave, sound intensity, or observed frequency.
    \item Write the governing equation before substituting.
    \item Check units: radians are dimensionless, but \(\omega\), \(k\), \(f\),
    and \(T\) are not interchangeable.
    \item For SHM, ask where the object is in its cycle: extreme, equilibrium,
    or between.
    \item For waves, ask whether you are tracking the wave shape or one element
    of the medium.
    \item For sound, decide whether the issue is speed, intensity, interference,
    pipe resonance, beats, Doppler shift, or shock waves.
\end{enumerate}

\section{Formula Sheet}

\begin{multicols}{2}

\subsection*{SHM}
\[
    x=x_m\cos(\omega t+\phi)
\]
\[
    v=-\omega x_m\sin(\omega t+\phi)
\]
\[
    a=-\omega^2x
\]
\[
    \omega=\frac{2\pi}{T}=2\pi f
\]

\subsection*{Spring Oscillator}
\[
    F=-kx
\]
\[
    \omega=\sqrt{\frac{k}{m}},
    \qquad
    T=2\pi\sqrt{\frac{m}{k}}
\]
\[
    E=\frac{1}{2}kx_m^2
\]
\[
    K=\frac{1}{2}k(x_m^2-x^2)
\]

\subsection*{Pendulums}
\[
    T=2\pi\sqrt{\frac{L}{g}}
\]
\[
    T=2\pi\sqrt{\frac{I}{mgh}}
\]

\subsection*{Waves}
\[
    y=y_m\sin(kx-\omega t)
\]
\[
    k=\frac{2\pi}{\lambda},
    \quad
    \omega=2\pi f,
    \quad
    v=\lambda f
\]
\[
    v=\sqrt{\frac{\tau}{\mu}}
\]
\[
    P_{\text{avg}}=\frac{1}{2}\mu v\omega^2y_m^2
\]

\subsection*{Standing Waves}
\[
    y'=2y_m\sin(kx)\cos(\omega t)
\]
\[
    f_n=\frac{nv}{2L}
\]

\subsection*{Sound}
\[
    v=\sqrt{\frac{B}{\rho}}
\]
\[
    I=\frac{P}{A},
    \qquad
    I=\frac{P_s}{4\pi r^2}
\]
\[
    \beta=(\SI{10}{dB})\log_{10}\left(\frac{I}{I_0}\right)
\]
\[
    f_{\text{beat}}=|f_1-f_2|
\]
\[
    f'=f\frac{v\pm v_D}{v\pm v_S}
\]

\end{multicols}

\section{Common Mistakes}

\begin{enumerate}
    \item Confusing frequency \(f\) with angular frequency \(\omega\).
    \item Saying speed is maximum at the turning points in SHM.
    \item Forgetting that \(a=-\omega^2x\), so acceleration points back toward
    equilibrium.
    \item Using amplitude when the problem asks for wavelength.
    \item Treating wave speed as the speed of one string element.
    \item Using \(kx+\omega t\) and \(kx-\omega t\) without checking direction.
    \item Assuming every standing-wave system has all harmonics; closed-open
    pipes have only odd harmonics.
    \item Adding decibels as though they were ordinary intensities.
    \item Choosing Doppler signs mechanically instead of checking whether the
    observed frequency should increase or decrease.
\end{enumerate}

\end{document}
